\section{Заключение}

В ходе выполнения курсовой работы было разработано приложение телефонного справочника с графическим интерфейсом. Все требования, указанные в постановке задачи, были успешно реализованы:

\begin{itemize}
    \item создание, редактирование и удаление контактов с полями: фамилия, имя, отчество, адрес, дата рождения, email, список телефонов с типами;
    \item валидация вводимых данных с использованием регулярных выражений: проверка ФИО на заглавную букву и допустимые символы, проверка формата телефона и email, проверка корректности даты рождения с учётом високосных лет;
    \item поиск контактов по различным полям с фильтрацией в реальном времени;
    \item сортировка таблицы по столбцам с трёхступенчатым переключением: по возрастанию, по убыванию, сброс к исходному порядку;
    \item сохранение данных в бинарный формат с использованием \texttt{QDataStream};
    \item экспорт и импорт контактов в формате YAML с использованием библиотеки yaml-cpp.
\end{itemize}

\noindentВ процессе разработки были изучены и освоены следующие технологии и концепции:

\begin{itemize}
    \item фреймворк Qt 6: архитектура приложений, механизм сигналов и слотов, система метаобъектов;
    \item виджеты Qt: \texttt{QMainWindow}, \texttt{QDialog}, \texttt{QTableWidget}, \texttt{QLineEdit}, \texttt{QDateEdit}, \texttt{QComboBox}, \texttt{QPushButton}, \texttt{QGroupBox};
    \item работа с файлами: класс \texttt{QFile}, бинарная сериализация через \texttt{QDataStream};
    \item регулярные выражения в Qt: класс \texttt{QRegularExpression} с поддержкой Unicode;
    \item наследование от базовых классов Qt и переопределение виртуальных методов;
    \item интеграция сторонних библиотек (yaml-cpp) через CMake и FetchContent;
    \item форматы сериализации данных: бинарный формат с магическим числом и версионированием, текстовый формат YAML.
\end{itemize}

\noindentДля разработки приложения использовались следующие инструменты:

\begin{itemize}
    \item операционная система: macOS 26.1;
    \item фреймворк: Qt 6;
    \item система сборки: CMake;
    \item среда разработки: Visual Studio Code версии 1.106.3;
    \item библиотека для работы с YAML: yaml-cpp.
\end{itemize}

Разработанное приложение демонстрирует практическое применение объектно-ориентированного программирования на C++ с использованием современного фреймворка Qt для создания кроссплатформенных приложений с графическим интерфейсом.